\chapter{Antidepressants}

\section*{Definition and Overview}
Antidepressants are medicines used primarily to treat moderate to severe depression, but they are also effective for a range of other conditions, including generalised anxiety disorder, panic disorder, obsessive-compulsive disorder, post-traumatic stress disorder, chronic pain, and some eating and sleep disorders. They are not stimulants: they do not lift mood within minutes. Instead, taken regularly over weeks, they gradually relieve the core symptoms of depression such as persistent low mood, loss of interest, poor sleep, low energy, and impaired concentration.

The main classes are \textbf{selective serotonin reuptake inhibitors} (SSRIs), such as fluoxetine, sertraline, escitalopram, and citalopram; \textbf{serotonin-noradrenaline reuptake inhibitors} (SNRIs), such as venlafaxine and duloxetine; and other agents including mirtazapine, bupropion, and the older \textbf{tricyclic antidepressants}. Antidepressants are most effective as part of a broader treatment plan that combines medication with psychological therapy, structured activity, and attention to sleep, exercise, and social support.

\section*{Epidemiology}
Depression is one of the most common health problems worldwide, affecting around one in six adults at some point in their lives and roughly one in ten in any given year. Antidepressants are among the most widely prescribed medicines in the world, and prescribing has risen steadily over recent decades in most high-income countries.

Women are prescribed antidepressants more often than men, reflecting both a higher reported prevalence of depression and anxiety and differences in help-seeking. A substantial proportion of prescriptions are written for conditions other than depression, notably anxiety disorders and chronic pain. Despite widespread use, a large number of people remain untreated or undertreated, and access to psychological therapy varies considerably between regions.

\section*{Aetiology and Pathophysiology}
The classic explanation that depression results simply from low levels of serotonin is now recognised as incomplete. Current understanding focuses on disrupted signalling between nerve cells, dysregulation of the stress response, and reduced \textbf{neuroplasticity} -- the brain's ability to remodel connections between neurons. Antidepressants work through several interacting mechanisms:

\begin{itemize}
    \item \textbf{Enhanced monoamine signalling:} SSRIs and SNRIs raise the concentration of serotonin, noradrenaline, or dopamine in the synapses between nerve cells by blocking their reuptake
    \item \textbf{Delayed neuroplasticity:} over several weeks of regular dosing, the brain adapts -- forming new synaptic connections and increasing growth factors such as BDNF -- which may explain why benefit takes 2--6 weeks to appear
    \item \textbf{Stress system modulation:} long-term stress and elevated cortisol are associated with depression and with changes in brain regions such as the hippocampus; antidepressants partially reverse these changes
    \item \textbf{Genes and environment:} heredity, childhood adversity, life events, and medical illness interact in complex, individual ways
\end{itemize}

\section*{Clinical Features}
When taking an antidepressant, the features observed relate to both the intended therapeutic effect and the side-effect profile:

\begin{itemize}
    \item \textbf{Gradual improvement:} lifting of mood, better sleep and appetite, more energy, and reduced anxiety over 2--6 weeks, sometimes preceded by a brief worsening
    \item \textbf{Early side effects:} nausea, headache, dry mouth, and a temporary increase in anxiety or restlessness in the first days to weeks
    \item \textbf{Ongoing effects:} sleep disturbance, sweating, weight change, and reduced sexual desire or function
    \item \textbf{Drowsiness:} more prominent with mirtazapine and some tricyclic antidepressants
    \item \textbf{Withdrawal symptoms:} dizziness, nausea, headache, and mood swings if the medicine is stopped abruptly or a dose is missed
\end{itemize}

Many early side effects settle within one to two weeks, which is one reason patients should be encouraged not to stop at the first sign of discomfort.

\section*{Differential Diagnosis}
When someone feels worse or develops new symptoms while taking an antidepressant, the possibilities to consider include:

\begin{itemize}
    \item \textbf{Worsening depression or anxiety} versus a \textbf{side effect} of the medicine
    \item \textbf{Withdrawal symptoms} after a missed or reduced dose versus a return of the underlying illness
    \item \textbf{Emerging mania} in bipolar disorder (elevated mood, racing thoughts, risky behaviour) versus a genuinely good response
    \item \textbf{Physical causes} of fatigue and low mood, such as hypothyroidism, anaemia, or vitamin B12 deficiency
    \item \textbf{Drug interactions} with other medicines, herbal remedies such as St John's wort, or alcohol
    \item \textbf{Pseudo-resistance} from poor adherence or an inadequate dose, rather than true treatment failure
\end{itemize}

\section*{When to Seek Medical Care}
See a doctor promptly if depression is severe, worsening, or has not improved after several weeks of treatment, or if side effects are intolerable. A review is also needed before any change in dose and before stopping the medicine.

\textbf{Call 000 (or local emergency services) for:}
\begin{itemize}
    \item Thoughts of suicide, self-harm, or plans to end your life
    \item Severe agitation, panic, or confusion early in treatment
    \item Confusion, fever, severe muscle stiffness or jerking, sweating, and rapid heart rate -- possible signs of \textbf{serotonin syndrome}
    \item Collapse, fainting, or any sign of overdose
\end{itemize}

\textbf{Immediate help is also available} from Lifeline (13 11 14) and other crisis support services, and from the nearest emergency department.

\section*{Diagnosis and Investigations}
\begin{itemize}
    \item \textbf{Clinical interview:} careful assessment of mood, sleep, appetite, energy, concentration, anxiety, and risk, along with past history and family history
    \item \textbf{Structured questionnaires:} tools such as the PHQ-9 and GAD-7 to measure severity, monitor progress, and guide dose decisions
    \item \textbf{Blood tests:} thyroid function, full blood count, and vitamin B12 or folate levels to exclude physical causes of low mood
    \item \textbf{Treatment review:} checking dose, adherence, duration, and whether an adequate trial of the chosen medicine has occurred
    \item \textbf{Monitoring for mania and risk:} watching for a switch to elevated mood in younger patients and repeated assessment of suicidal risk, especially in the first weeks and after dose increases
    \item \textbf{Baseline ECG and electrolytes:} considered with citalopram or tricyclics, which can affect heart rhythm, and in older patients
\end{itemize}

\section*{Management}
\begin{itemize}
    \item \textbf{Take as prescribed:} daily, at a consistent time, for weeks or months; benefit usually appears over 2--6 weeks
    \item \textbf{Give an adequate trial:} a dose may need to be increased and a full 4--8 week trial allowed before judging effectiveness
    \item \textbf{Switch or combine:} if there is no response, options include switching class, adding a second agent, or specialist review
    \item \textbf{Continue after recovery:} most people remain on the medicine for 6--12 months after remission to prevent relapse, and longer for recurrent depression
    \item \textbf{Psychological therapy:} cognitive behaviour therapy, and for severe cases behavioural activation and mindfulness-based approaches, improve outcomes when combined with medication
    \item \textbf{Lifestyle measures:} regular exercise, sleep hygiene, and reduced alcohol use support recovery
    \item \textbf{Tapering:} stopping is always done gradually under medical supervision to minimise withdrawal symptoms
    \item \textbf{Specialist referral:} for treatment-resistant depression, bipolar disorder, severe or psychotic depression, or pregnancy planning
\end{itemize}

\section*{Complications}
\begin{itemize}
    \item Serotonin syndrome, particularly when serotonergic medicines, opioids, or St John's wort are combined
    \item Hyponatraemia (low sodium), mainly in older people in the first weeks of an SSRI
    \item Increased bleeding risk, especially when combined with NSAIDs or anticoagulants
    \item Sexual dysfunction and weight gain, which commonly affect adherence
    \item Withdrawal symptoms and rebound relapse when medicines are stopped abruptly
    \item A small rise in suicidal ideation in young people early in treatment, requiring close monitoring
    \item Overdose toxicity, especially with tricyclic antidepressants, which can cause dangerous heart rhythm disturbances
\end{itemize}

\section*{Prognosis and Outcomes}
Most people with depression improve with treatment; around half to two-thirds respond to their first antidepressant, and a substantial further proportion respond when the medicine is changed or combined with therapy. Full remission takes time and may require dose adjustment or a switch. Continuing the medicine after recovery substantially reduces the chance of relapse, and for people with recurrent depression long-term treatment is often worthwhile.

Outcomes are best when medication is combined with psychological therapy, social support, and healthy routines, and when the person is actively involved in decisions about treatment. Antidepressants are not addictive in the usual sense, but they should be withdrawn gradually; for most people, the benefits far outweigh the risks when the medicine is used appropriately.

\section*{Prevention and Self-Care}
\begin{itemize}
    \item Maintain regular exercise, a consistent sleep routine, and a balanced diet
    \item Limit alcohol, which is a depressant and interacts with these medicines
    \item Stay connected with supportive people, and seek help early when mood dips
    \item Take medicines reliably and never stop them without medical advice
    \item Learn the early warning signs of relapse and agree on a plan with your doctor
    \item Report side effects, including sexual problems, rather than quietly stopping the medicine
    \item Keep all medicines and supplements out of reach of children and store them safely
\end{itemize}

\section*{Sources and Further Reading}
The following organisations provide reliable, up-to-date information on antidepressants, depression, and treatment for patients, students, and health professionals.

\begin{itemize}
    \item NHS. \textit{Information on antidepressants and treatment for depression.} https://www.nhs.uk/conditions/
    \item NICE. \textit{Guidelines on the treatment and management of depression.} https://www.nice.org.uk/
    \item Mayo Clinic. \textit{Overview of antidepressants, uses, and side effects.} https://www.mayoclinic.org/
    \item MedlinePlus. \textit{Health information on antidepressants and drug interactions.} https://medlineplus.gov/
    \item World Health Organization. \textit{Resources on mental health and depression.} https://www.who.int/
    \item MSD Manual. \textit{Professional overview of depression and drug treatment.} https://www.msdmanuals.com/
\end{itemize}
